\makeatletter
\def\input@path{{./}{../}{../../}{../../../}{../../../../}{preamble/}{../preamble/}{../../preamble/}{../../../preamble/}{../../../../book2026-2027/preamble/}}
\makeatother
\input{preamble_exams.tex}
\input{graphs/highlights.tex}
% \GTenHighlightsfalse % Uncomment to disable instructional highlights.
% Population/model, sample evidence, design threshold, and result roles.
\newcommand{\Pop}[1]{\HighlightStroke{GTenHLParameter}{#1}}
\newcommand{\Obs}[1]{\HighlightStroke{GTenHLValue}{#1}}
\newcommand{\Design}[1]{\HighlightStroke{GTenHLThird}{#1}}
\newcommand{\Result}[1]{\HighlightStroke{GTenHLOutcome}{#1}}

\title{Grade 11 Global Economics\\Term 3 --- Lesson 1\\Home Practice C3: Type I and Type II errors}
\author{}\date{}
\hypersetup{pdftitle={Grade 11 Global Economics - Term 3 - Lesson 1 - C3},pdfsubject={Type I and Type II errors}}
\begin{document}
\maketitle
\setcounter{tocdepth}{1}
\setcounter{secnumdepth}{2}
\phantomsection\label{toc}
\tableofcontents
\section*{Evaluated criteria}
\begin{description}
\item[\textbf{C3}] Distinguishes Type I and Type II errors by identifying the incorrect decision associated with each and explaining their consequences in economic and financial contexts.
\end{description}
\textbf{How to practise.} Write decision and possible population truth separately, then describe the economic consequence. Attempt each task before opening its solution link. Use the return arrows to revisit the problem, group, or contents. All datasets and scenarios in this activity are illustrative.
\clearpage
\section{Introduction}\label{sec:introduction}
\subsection*{Highlight legend}
\phantomsection\label{tab:highlight-legend}
\begin{tabularx}{\linewidth}{@{}>{\raggedright\arraybackslash}p{.22\linewidth}>{\raggedright\arraybackslash}p{.25\linewidth}>{\raggedright\arraybackslash}X@{}}
\toprule
\textbf{Symbol or notation}&\textbf{Concept}&\textbf{Brief meaning}\\\midrule
$\Pop{H_0},\Result{H_a}$ & \Pop{Claims} & Population statements. \\
$\Result{\mathrm{I}}$ & \Result{Type I error} & Reject a true null. \\
$\Result{\mathrm{II}}$ & \Result{Type II error} & Fail to reject a false null. \\
$\Design{\alpha}$ & \Design{Level} & False-rejection control. \\
\bottomrule
\end{tabularx}
\subsection{Introductory example 1: Fuel budget}\label{ex:intro-one}
A transport authority tests $H_0:\mu=4.0$ against $H_a:\mu>4.0$ for mean monthly fuel inflation in percentage points at $\Design{\alpha}=\Design{0.05}$. Rejection when the true mean is 4.0\% is Type I; fares might rise unnecessarily.
\subsection{Introductory example 2: Employment reform}\label{ex:intro-two}
A ministry tests $H_0:\mu=11$ against $H_a:\mu<11$ for mean unemployment duration in weeks. Non-rejection when the true mean is 9 weeks is Type II; a beneficial reform might lose funding.
\section{Notation Reference and Practice Procedure}\label{sec:notation-and-procedure}
\subsection{Notation reference}\label{tab:notation-reference}
\begin{tabularx}{\linewidth}{@{}>{\raggedright\arraybackslash}p{.16\linewidth}>{\raggedright\arraybackslash}p{.18\linewidth}>{\raggedright\arraybackslash}p{.28\linewidth}>{\raggedright\arraybackslash}X@{}}
\toprule
\textbf{Symbol or notation}&\textbf{Concept}&\textbf{Definition}&\textbf{Meaning in the introductory examples}\\\midrule
$\Pop{H_0},\Result{H_a}$ & \Pop{Claims} & Truth concerns a population; decisions use sample evidence. & \textbullet\ \hyperref[ex:intro-one]{Example 1}: fuel inflation.\par\textbullet\ \hyperref[ex:intro-two]{Example 2}: unemployment duration. \\
$\Result{\mathrm{I}}$ & \Result{Type I error} & False rejection requires both rejection and a true null. & \textbullet\ \hyperref[ex:intro-one]{Example 1}: flag an increase at a true mean of 4.0\%. \\
$\Result{\mathrm{II}}$ & \Result{Type II error} & A missed population departure in the alternative direction. & \textbullet\ \hyperref[ex:intro-two]{Example 2}: miss a reduction below 11 weeks. \\
$\Design{\alpha}$ & \Design{Level} & Rejection probability under the null, not the probability an observed conclusion is wrong. & \textbullet\ \hyperref[ex:intro-one]{Example 1}: 0.05. \\
\bottomrule
\end{tabularx}
\subsection{Practice procedure}
\begin{enumerate}
\item Define the population parameter and state the supplied hypotheses.
\item Identify the decision separately from any assumed truth.
\item Place decision and truth in the two-by-two framework.
\item Describe the false claim or missed effect, who acts, and the economic consequence.
\item When truth is unknown, describe a possible error rather than asserting that one occurred.
\end{enumerate}
\section{General Worked Problems}\label{sec:general-worked}
\SectionProblemsTableOfContents{\SectionProblemLink{prob:matrix}{Decision versus truth} \\
\SectionProblemLink{prob:reports}{Two policy reports}}
\subsection{Decision versus truth}\label{prob:matrix}
A bank tests $H_0:p=0.05$ against $H_a:p>0.05$ for the one-year default proportion. Rejection triggers a costly lending review. Assume a valid sampling model; no calculation is required.
\QuestionsTasks[sec:general-worked]{{3}/{sol:matrix}/{Classify all four decision-versus-truth combinations and explain both error costs.}}
\subsubsection*{Worked solution}\label{sol:matrix}
\begin{center}\begin{tabularx}{\linewidth}{lXX}\toprule Decision & $H_0$ true & $H_0$ false in the alternative direction\\\midrule Reject $H_0$ & \Result{Type I}: false alarm & Correct detection\\ Fail to reject $H_0$ & Correct non-rejection & \Result{Type II}: missed excess risk\\\bottomrule\end{tabularx}\end{center}
A Type I error announces excess defaults when the true rate is 5\%; unnecessary reviews and tighter credit cost staff time and may exclude reliable borrowers. Type II misses a rate above 5\%; underpriced lending may continue and cause losses. The same observed decision can be correct or incorrect depending on truth. A test decision does not reveal which truth column applies.
\secInicio[prob:matrix][sec:general-worked]
\subsection{Two policy reports}\label{prob:reports}
For the introductory fuel hypotheses, a valid report gives p-value 0.03 at $\alpha=0.05$. For the employment hypotheses, a valid report gives p-value 0.18 at $\alpha=0.10$. The sampling and distribution conditions are satisfied.
\QuestionsTasks[sec:general-worked]{{3}/{sol:reports}/{State decisions and possible errors with consequences. Correct the claim that the fuel null has a 3 percent chance of being true.}}
\subsubsection*{Worked solution}\label{sol:reports}
Fuel: $0.03<0.05$, so \Result{reject $H_0$}; evidence supports population mean inflation above 4.0\%. A possible Type I error is announcing an increase at a true mean of 4.0\%, causing unnecessary fare increases. Employment: $0.18>0.10$, so \Result{fail to reject $H_0$}; evidence is insufficient for a mean below 11 weeks. A possible Type II error misses a real reduction and may withdraw useful support. Neither error is known to have happened. The fuel p-value is the probability of results at least as extreme under the null model, not the probability the null is true. Non-rejection does not establish no effect.
\secInicio[prob:reports][sec:general-worked]
\section{Advanced Problems}\label{sec:advanced}
\SectionProblemsTableOfContents{\SectionProblemLink{prob:costs}{Error costs and investment decisions}}
\subsection{Error costs and investment decisions}\label{prob:costs}
A fund tests $H_0:\mu=0$ versus $H_a:\mu>0$ for mean net excess return after fees. Rejection triggers a 2 million USD expansion; non-rejection leaves the pilot unchanged. In A, the fund rejects and true mean is zero. In B, it fails to reject and true mean is positive. A manager says the more expensive action must always be Type II.
\QuestionsTasks[sec:advanced]{{3}/{sol:costs}/{Classify A and B, assess the manager's rule, and explain what could be said if truth were unknown.}}
\subsubsection*{Worked solution}\label{sol:costs}
A is \Result{Type I}: rejection of a true no-excess-return benchmark. Expansion can create unnecessary operating costs and uncompensated risk. B is \Result{Type II}: failure to detect a genuine positive mean, possibly losing profitable opportunities. Error names depend on decision and truth, not the budget size or whether action sounds cautious. Costs help select a design but do not redefine errors. With unknown truth, rejection carries a possible Type I error and non-rejection a possible Type II error. Neither proves truth. Even a genuine positive mean might be too small to cover expansion costs; statistical evidence and practical value are separate.
\secInicio[prob:costs][sec:advanced]
\end{document}

